\chapter{Derivación: un primer curso}\label{ch:g11:deriv}

¿Con qué rapidez cambia una \omterm{def:g11:func:function}{función}? La respuesta en un punto concreto es
la \emph{\omterm{def:g11:deriv:derivative}{derivada}}, el número más útil de todas las matemáticas aplicadas.
Este capítulo la construye a partir de las \omterm{def:g11:deriv:rate}{tasas de variación medias}, la
calcula para las \omterm{def:g11:func:reference}{funciones de referencia} y usa su signo para leer la
variación de una \omterm{def:g11:func:function}{función}. Todo esto se amplía y se profundiza en el
\cref{ch:g12:deriv}.

\section{La tasa de variación media}

\begin{definition}[Tasa de variación media]\label{def:g11:deriv:rate}
Sea $f$ definida en un \omterm{def:g10:numbers:interval}{intervalo} $I$ y sean $a \neq b$ en $I$. La
\emph{tasa de variación media}\index{tasa de variación media} de $f$ entre
$a$ y $b$ es
\[
\frac{f(b) - f(a)}{b - a}.
\]
Es la \omterm{def:g10:reffunc:affine}{pendiente} de la recta \emph{secante}\index{secante} que pasa por los
puntos $\bigl(a, f(a)\bigr)$ y $\bigl(b, f(b)\bigr)$ de la curva.
\end{definition}

\begin{example}
Un coche recorre $140$ km entre las 2 y las 4 de la tarde: su velocidad
\omterm{def:g10:stats:mean}{media} es $\frac{140}{2} = 70$ km/h. El velocímetro, en cambio, muestra una
velocidad \emph{instantánea} en cada momento; esa es exactamente la noción
que la \omterm{def:g11:deriv:derivative}{derivada} hace precisa.
\end{example}

\section{La derivada en un punto}

\begin{definition}[Derivada en un punto]\label{def:g11:deriv:derivative}
Sea $f$ definida en un \omterm{def:g10:numbers:interval}{intervalo} que contiene a $a$. Para $h \neq 0$,
formamos la \omterm{def:g11:deriv:rate}{tasa de variación media} entre $a$ y $a + h$:
\[
r(h) = \frac{f(a+h) - f(a)}{h}.
\]
Si $r(h)$ se acerca a un único número fijo cuando $h$ se aproxima
arbitrariamente a $0$, decimos que $f$ es \emph{derivable en
$a$}\index{derivada}\index{derivable}; ese número es la \emph{derivada de
$f$ en $a$}, que se escribe $f'(a)$:
\[
f'(a) = \lim_{h \to 0} \frac{f(a+h) - f(a)}{h}.
\]
\end{definition}

\begin{example}\label{ex:g11:deriv:atpoint}
Sea $f(x) = x^2$ y $a = 1$. Para $h \neq 0$:
\[
\frac{f(1+h) - f(1)}{h}
= \frac{(1+h)^2 - 1}{h}
= \frac{1 + 2h + h^2 - 1}{h}
= \frac{h(2 + h)}{h}
= 2 + h .
\]
Cuando $h$ se acerca a $0$, esto se acerca a $2$: la \omterm{def:g11:func:function}{función} $x^2$ es
\omterm{def:g11:deriv:derivative}{derivable} en $1$ y $f'(1) = 2$. Fíjate en la estrategia: \emph{simplificar
el cociente hasta que $h$ deje de aparecer en un denominador} y hacer
después $h \to 0$.
\end{example}

\begin{example}[Un punto sin derivada]
Sea $f(x) = \abs{x}$ y $a = 0$. Entonces
$\frac{f(h) - f(0)}{h} = \frac{\abs h}{h}$, que vale $1$ para $h > 0$ y
$-1$ para $h < 0$: no se acerca a ningún número único, así que $\abs x$ no
es \omterm{def:g11:deriv:derivative}{derivable} en $0$. Su curva tiene ahí un \emph{pico}.
\end{example}

\section{La recta tangente}

\begin{definition}[Tangente]\label{def:g11:deriv:tangent}
Si $f$ es \omterm{def:g11:deriv:derivative}{derivable} en $a$, la \emph{tangente}\index{tangente} a la curva
en el punto $A = \bigl(a, f(a)\bigr)$ es la recta que pasa por $A$ con
\omterm{def:g10:reffunc:affine}{pendiente} $f'(a)$. Su \omterm{def:g10:algebra:equation}{ecuación} es
\[
y = f(a) + f'(a)(x - a).
\]
\end{definition}

La \omterm{def:g11:deriv:tangent}{tangente} es la posición límite de las \omterm{def:g11:deriv:rate}{secantes} que pasan por $A$:
cuando $h \to 0$, el segundo punto $\bigl(a+h, f(a+h)\bigr)$ se desliza
por la curva hacia $A$ y la \omterm{def:g10:reffunc:affine}{pendiente} de la \omterm{def:g11:deriv:rate}{secante},
$\frac{f(a+h)-f(a)}{h}$, se acerca a $f'(a)$.

\begin{omfigure}
\begin{tikzpicture}
\begin{axis}[omaxis,
  width=8.8cm, height=6cm,
  xmin=-0.4, xmax=3.2, ymin=-0.6, ymax=4.3,
  xtick={1,2.2}, xticklabels={$a$,$a+h$}, ytick=\empty]
  \addplot[omProp, thick, domain=-0.1:2.9] {1 + 1.6*(x-1)};
  \addplot[omThm, thick, domain=-0.35:2.6] {1 + 1*(x-1)};
  \addplot[omDef, thick, domain=-0.4:2.75] {x^2/2 + 0.5};
  \fill (axis cs:1,1) circle (1.5pt) node[above left, font=\small] {$A$};
  \fill (axis cs:2.2,2.92) circle (1.5pt)
    node[above left, font=\small] {$B$};
  \node[omProp, font=\small, anchor=west] at (axis cs:2.62,3.75)
    {secante};
  \node[omThm, font=\small, anchor=west] at (axis cs:2.62,2.45)
    {tangente};
\end{axis}
\end{tikzpicture}

{\small La \omterm{def:g11:deriv:rate}{secante} que pasa por $A$ y $B = \bigl(a+h, f(a+h)\bigr)$ tiene
\omterm{def:g10:reffunc:affine}{pendiente} $\frac{f(a+h)-f(a)}{h}$. Cuando $B$ se desliza hacia $A$, la
\omterm{def:g11:deriv:rate}{secante} se inclina hasta convertirse en la \omterm{def:g11:deriv:tangent}{tangente}, cuya \omterm{def:g10:reffunc:affine}{pendiente} es
$f'(a)$.}
\end{omfigure}

\begin{example}\label{ex:g11:deriv:tangenteq}
Para $f(x) = x^2$ en $a = 1$: $f(1) = 1$ y $f'(1) = 2$
(\cref{ex:g11:deriv:atpoint}), así que la \omterm{def:g11:deriv:tangent}{tangente} es
\[
y = 1 + 2(x - 1) = 2x - 1 .
\]
Es la recta $d_2$ hallada por álgebra pura en el
\cref{exo:g11:quad:11}.
\end{example}

\section{Derivadas de las funciones de referencia}

\begin{proposition}\label{prop:g11:deriv:refderivatives}
En sus \omterm{def:g11:func:function}{dominios} de derivabilidad:
\[
(c)' = 0, \quad
(x)' = 1, \quad
(x^2)' = 2x, \quad
(x^3)' = 3x^2, \quad
\left(\frac1x\right)' = -\frac{1}{x^2}, \quad
(\sqrt x\,)' = \frac{1}{2\sqrt x} \ (x > 0).
\]
Más en general, $(x^n)' = n\,x^{n-1}$ para todo entero $n \geq 1$.
\end{proposition}

\begin{proof}
Cada demostración sigue la estrategia del \cref{ex:g11:deriv:atpoint}:
simplificar la tasa y hacer después $h \to 0$. Fijemos $a$ en el \omterm{def:g11:func:function}{dominio}.

\emph{Constante e identidad.} $\frac{c - c}{h} = 0$ y
$\frac{(a+h) - a}{h} = 1$ para todo $h$: los límites son $0$ y $1$.

\emph{Cuadrado.} $\frac{(a+h)^2 - a^2}{h} = \frac{2ah + h^2}{h} = 2a + h
\to 2a$.

\emph{Cubo.} $(a+h)^3 = a^3 + 3a^2h + 3ah^2 + h^3$, así que la tasa es
$3a^2 + 3ah + h^2 \to 3a^2$.

\emph{Inverso.} Para $a \neq 0$ y $h$ suficientemente pequeño para que
$a + h \neq 0$:
\[
\frac{1}{h}\left(\frac{1}{a+h} - \frac{1}{a}\right)
= \frac{1}{h} \cdot \frac{a - (a+h)}{a(a+h)}
= \frac{-1}{a(a+h)}
\xrightarrow[h \to 0]{} -\frac{1}{a^2}.
\]

\emph{\omterm{def:g11:quad:discriminant}{Raíz} cuadrada.} Para $a > 0$, multiplicamos por el conjugado:
\[
\frac{\sqrt{a+h} - \sqrt a}{h}
= \frac{(a + h) - a}{h\,(\sqrt{a+h} + \sqrt a)}
= \frac{1}{\sqrt{a+h} + \sqrt a}
\xrightarrow[h \to 0]{} \frac{1}{2\sqrt a}.
\]

\emph{Potencia general.} El patrón $2a$, $3a^2$ continúa: al \omterm{def:g10:algebra:expand}{desarrollar}
$(a+h)^n$, todos los términos posteriores a $a^n + n\,a^{n-1}h$ llevan
$h^2$, así que la tasa tiende a $n\,a^{n-1}$. (En el \cref{ch:g12:deriv}
se da una inducción completa.)
\end{proof}

\section{Reglas de derivación}

\begin{theorem}[Operaciones]\label{thm:g11:deriv:rules}
Sean $u$ y $v$ \omterm{def:g11:deriv:derivative}{derivables} en $I$ y sea $\lambda \in \R$. Entonces
$u + v$, $\lambda u$ y $uv$ son \omterm{def:g11:deriv:derivative}{derivables} en $I$, igual que
$\frac{u}{v}$ allí donde $v \neq 0$, y
\[
(u+v)' = u' + v', \qquad
(\lambda u)' = \lambda u', \qquad
(uv)' = u'v + uv', \qquad
\left(\frac{u}{v}\right)' = \frac{u'v - uv'}{v^2}.
\]
\end{theorem}

\begin{proof}[Demostración de las tres primeras reglas]
Fijemos $a \in I$ y abreviemos $\Delta u = u(a+h) - u(a)$ y
$\Delta v = v(a+h) - v(a)$.

\emph{Suma.} La tasa de $u + v$ es
$\frac{\Delta u + \Delta v}{h} = \frac{\Delta u}{h} + \frac{\Delta v}{h}
\to u'(a) + v'(a)$.

\emph{Múltiplo.} La tasa de $\lambda u$ es
$\lambda \frac{\Delta u}{h} \to \lambda u'(a)$.

\emph{Producto.} Sumamos y restamos el término mixto $u(a+h)\,v(a)$:
\begin{align*}
\frac{u(a+h)v(a+h) - u(a)v(a)}{h}
&= \frac{u(a+h)v(a+h) - u(a+h)v(a) + u(a+h)v(a) - u(a)v(a)}{h}\\
&= u(a+h)\,\frac{\Delta v}{h} + v(a)\,\frac{\Delta u}{h}.
\end{align*}
Cuando $h \to 0$: $\frac{\Delta v}{h} \to v'(a)$,
$\frac{\Delta u}{h} \to u'(a)$ y $u(a+h) \to u(a)$ (su tasa tiene límite,
luego $\Delta u = h \cdot \frac{\Delta u}{h} \to 0$). Por tanto, la tasa
tiende a $u(a)v'(a) + v(a)u'(a)$.

La regla del cociente se demuestra por el mismo camino; aquí se
\emph{admite a este nivel} y se demuestra en el \cref{ch:g12:deriv}.
\end{proof}

\begin{example}\label{ex:g11:deriv:computations}
\emph{Polinomio.} $f(x) = 3x^3 - 5x^2 + 7x - 2$:
\[
f'(x) = 3 \times 3x^2 - 5 \times 2x + 7 = 9x^2 - 10x + 7 .
\]

\emph{Cociente.} $g(x) = \frac{2x+1}{x^2+1}$: con $u = 2x + 1$ y
$v = x^2 + 1$,
\[
g'(x) = \frac{2(x^2+1) - (2x+1)\,2x}{(x^2+1)^2}
= \frac{-2x^2 - 2x + 2}{(x^2+1)^2}.
\]

\emph{Producto.} $k(x) = x\sqrt x$ para $x > 0$:
$k'(x) = 1 \times \sqrt x + x \times \frac{1}{2\sqrt x}
= \sqrt x + \frac{\sqrt x}{2} = \frac{3\sqrt x}{2}$.
\end{example}

\section{Signo de la derivada y variación}

\begin{theorem}[Derivada y variación]\label{thm:g11:deriv:variations}
Sea $f$ \omterm{def:g11:deriv:derivative}{derivable} en un \omterm{def:g10:numbers:interval}{intervalo} $I$.
\begin{enumerate}
  \item Si $f' \geq 0$ en $I$, entonces $f$ es \omterm{def:g11:func:monotone}{creciente} en $I$.
  \item Si $f' \leq 0$ en $I$, entonces $f$ es \omterm{def:g10:functions:variations}{decreciente} en $I$.
  \item Si $f' > 0$ en $I$ (salvo, quizá, en un número finito de puntos
        donde se anula), entonces $f$ es estrictamente \omterm{def:g11:func:monotone}{creciente} en $I$.
\end{enumerate}
\end{theorem}

\begin{proof}
La intuición está clara (una curva cuyas \omterm{def:g11:deriv:tangent}{tangentes} apuntan todas hacia
arriba tiene que subir), pero una demostración rigurosa necesita el
teorema del valor medio. El resultado se \emph{admite a este nivel};
véase el \cref{thm:g12:deriv:variations}.
\end{proof}

\begin{proposition}[Extremos]\label{prop:g11:deriv:extrema}
Si $f'$ cambia de signo en $a$ de $+$ a $-$, entonces $f$ tiene un \omterm{def:g10:functions:extrema}{máximo}
relativo en $a$; si cambia de $-$ a $+$, un \omterm{def:g10:functions:extrema}{mínimo} relativo. En un punto
así, $f'(a) = 0$ y la \omterm{def:g11:deriv:tangent}{tangente} es horizontal.
\end{proposition}

\begin{proof}
Si $f' \geq 0$ en un \omterm{def:g10:numbers:interval}{intervalo} que termina en $a$ y $f' \leq 0$ en un
\omterm{def:g10:numbers:interval}{intervalo} que empieza en $a$, entonces $f$ crece hasta $f(a)$ y decrece
después (\cref{thm:g11:deriv:variations}): $f(a)$ es el mayor valor de los
alrededores. Como $f'$ cambia de signo en $a$ y está definida allí,
$f'(a) = 0$.
\end{proof}

\begin{method}[Estudiar una función]\label{met:g11:deriv:study}
\begin{enumerate}
  \item Calcula $f'(x)$ y simplifícala, a ser posible hasta una forma
        factorizada.
  \item Determina el signo de $f'(x)$ en cada \omterm{def:g10:numbers:interval}{intervalo} (si $f'$ es de
        segundo grado, usa el \cref{thm:g11:quad:sign}).
  \item Recoge los resultados en una \emph{\omterm{def:g10:functions:table}{tabla de variación}}: una fila
        para el signo de $f'$ y otra de flechas para $f$, con los valores
        de $f$ en los puntos donde cambia el sentido.
  \item Concluye: extremos, número de soluciones de $f(x) = k$,
        desigualdades.
\end{enumerate}
\end{method}

\begin{example}\label{ex:g11:deriv:study}
Estudiemos $f(x) = x^3 - 3x + 1$ en $\R$. En primer lugar,
$f'(x) = 3x^2 - 3 = 3(x-1)(x+1)$: un trinomio de raíces $-1$ y $1$,
positivo fuera de ellas. Así pues, $f$ es \omterm{def:g11:func:monotone}{creciente} en
$\intoc{-\infty}{-1}$, \omterm{def:g10:functions:variations}{decreciente} en $\intcc{-1}{1}$ y \omterm{def:g11:func:monotone}{creciente} en
$\intco{1}{+\infty}$, con un \omterm{def:g10:functions:extrema}{máximo} relativo $f(-1) = -1 + 3 + 1 = 3$ y un
\omterm{def:g10:functions:extrema}{mínimo} relativo $f(1) = 1 - 3 + 1 = -1$.
\end{example}

\begin{omfigure}
\begin{tikzpicture}
\begin{axis}[omaxis,
  width=8.4cm, height=5.8cm,
  xmin=-2.6, xmax=2.6, ymin=-3.6, ymax=4.4,
  xtick={-1,1}, ytick={-1,3}]
  \addplot[omProp, thick, domain=-1.7:-0.3] {3};
  \addplot[omProp, thick, domain=0.3:1.7] {-1};
  \addplot[omDef, thick, domain=-2.25:2.25, samples=90] {x^3 - 3*x + 1};
  \fill (axis cs:-1,3) circle (1.5pt);
  \fill (axis cs:1,-1) circle (1.5pt);
\end{axis}
\end{tikzpicture}

{\small La curva de $f(x) = x^3 - 3x + 1$ (\cref{ex:g11:deriv:study}):
\omterm{def:g11:deriv:tangent}{tangentes} horizontales (en naranja) en el \omterm{def:g10:functions:extrema}{máximo} relativo $(-1, 3)$ y en
el \omterm{def:g10:functions:extrema}{mínimo} relativo $(1, -1)$, donde $f'$ cambia de signo.}
\end{omfigure}

\begin{example}[Optimización]\label{ex:g11:deriv:box}
Se fabrica una caja abierta a partir de una lámina de $12 \times 12$ cm
recortando un cuadrado de lado $x$ en cada esquina y doblando. Su volumen
es $V(x) = x(12 - 2x)^2$ para $0 < x < 6$. Desarrollando,
$V(x) = 4x^3 - 48x^2 + 144x$, luego
\[
V'(x) = 12x^2 - 96x + 144 = 12(x^2 - 8x + 12) = 12(x - 2)(x - 6).
\]
En $\intoo{0}{6}$, el factor $x - 6$ es negativo, así que $V'(x)$ tiene el
signo de $-(x-2)$: positivo antes de $2$ y negativo después. El volumen es
\omterm{def:g10:functions:extrema}{máximo} para $x = 2$ cm, y vale $V(2) = 2 \times 8^2 = 128$ cm$^3$.
\end{example}

\section{Ejercicios}

\begin{exercise}[$\star$]\label{exo:g11:deriv:1}
Usando la definición (la tasa de variación y después $h \to 0$), calcula
$f'(2)$ para $f(x) = x^2$ y después $g'(1)$ para $g(x) = x^2 + 3x$.
\end{exercise}

\begin{exercise}[$\star$]\label{exo:g11:deriv:2}
Deriva:
\[
f(x) = 4x^3 - 2x^2 + x - 7, \qquad
g(x) = 5\sqrt{x} + \frac{3}{x}, \qquad
h(x) = (2x+1)(x^2 - 3) .
\]
\end{exercise}

\begin{exercise}[$\star$]\label{exo:g11:deriv:3}
Da la \omterm{def:g10:algebra:equation}{ecuación} de la \omterm{def:g11:deriv:tangent}{tangente} a la curva de $f$ en el punto de \omterm{def:g10:coordgeom:system}{abscisa}
$a$:
\[
f(x) = x^2 - 3x, \ a = 2; \qquad
f(x) = \frac1x, \ a = 1; \qquad
f(x) = \sqrt x, \ a = 4 .
\]
\end{exercise}

\begin{exercise}[$\star$]\label{exo:g11:deriv:4}
Deriva los cocientes:
\[
f(x) = \frac{x+2}{x-1}, \qquad
g(x) = \frac{x^2}{x^2+1}, \qquad
h(x) = \frac{1}{x^2 + x + 1} .
\]
\end{exercise}

\begin{exercise}[$\star\star$]\label{exo:g11:deriv:5}
Estudia la variación de $f(x) = x^3 - 6x^2 + 9x - 2$ en $\R$ (\omterm{def:g10:functions:table}{tabla de
variación} y extremos relativos).
\end{exercise}

\begin{exercise}[$\star\star$]\label{exo:g11:deriv:6}
Estudia la variación de $f(x) = \frac{x^2 + 3}{x + 1}$ en
$\intoo{-1}{+\infty}$.
\end{exercise}

\begin{exercise}[$\star\star$]\label{exo:g11:deriv:7}
Sea $f(x) = x^2 - 4x + 5$.
\begin{enumerate}
  \item Halla el punto de la curva donde la \omterm{def:g11:deriv:tangent}{tangente} es horizontal.
  \item Halla el punto o los puntos donde la \omterm{def:g11:deriv:tangent}{tangente} es paralela a la
        recta $y = 2x + 1$.
\end{enumerate}
\end{exercise}

\begin{exercise}[$\star\star$]\label{exo:g11:deriv:8}
Se construye un cercado rectangular apoyado en un muro con $60$ m de valla
(tres lados). Expresa el área en \omterm{def:g11:func:function}{función} de la anchura $x$ y usa la
\omterm{def:g11:deriv:derivative}{derivada} para hallar las dimensiones de área máxima. (Compara con el
método del vértice del \cref{ch:g11:quad}.)
\end{exercise}

\begin{exercise}[$\star\star$]\label{exo:g11:deriv:9}
Demuestra que para todo $x > 0$ se cumple $\sqrt{x} \leq \frac{x + 1}{2}$,
estudiando la \omterm{def:g11:func:function}{función} $f(x) = \frac{x+1}{2} - \sqrt x$ en
$\intoo{0}{+\infty}$.
\end{exercise}

\begin{exercise}[$\star\star$]\label{exo:g11:deriv:10}
Una lata cilíndrica debe contener $500$ cm$^3$. Su superficie (los dos
discos más el lateral) es $S = 2\pi r^2 + 2\pi r h$ con $\pi r^2 h = 500$.
Expresa $S$ en \omterm{def:g11:func:function}{función} de $r$ sola y demuestra que
$S'(r) = 4\pi r - \frac{1000}{r^2}$. Deduce el radio que minimiza la
superficie, en forma exacta.
\end{exercise}

\begin{exercise}[$\star\star\star$]\label{exo:g11:deriv:11}
¿Cuántas \omterm{def:g11:deriv:tangent}{tangentes} a la \omterm{def:g11:quad:parabola}{parábola} $y = x^2$ pasan por el punto exterior
$P(1, -3)$? Determina sus ecuaciones. (Llama $a$ a la \omterm{def:g10:coordgeom:system}{abscisa} del punto de
contacto y expresa que la \omterm{def:g11:deriv:tangent}{tangente} en $a$ pasa por $P$.)
\end{exercise}

\begin{exercise}[$\star\star\star$]\label{exo:g11:deriv:12}
Sea $f(x) = x^3 - 3x + 1$ (véase el \cref{ex:g11:deriv:study}). Usando la
\omterm{def:g10:functions:table}{tabla de variación}, determina el número de soluciones de la \omterm{def:g10:algebra:equation}{ecuación}
$f(x) = k$ según el valor del \omterm{def:g10:numbers:sets}{número real} $k$.
\end{exercise}

\section{Problema: Los tres oficios de la tangente}

\begin{problem}[{Problema de fin de semana --- el espejo parabólico
demostrado, la máquina de Newton para hallar raíces y la viga más
resistente que da un tronco}]
\label{pb:g11:deriv:1}
Una \omterm{def:g11:deriv:derivative}{derivada} traza \omterm{def:g11:deriv:tangent}{tangentes}; suena a talento modesto. Este problema
muestra a la \omterm{def:g11:deriv:tangent}{tangente} desempeñando tres oficios a la vez: demuestra el
teorema del faro que quedó \omterm{def:g10:reffunc:affine}{pendiente} en el \cref{pb:g11:quad:1} (todo rayo
paralelo al eje de una \omterm{def:g11:quad:parabola}{parábola} se refleja pasando por el \omterm{pb:g11:quad:1}{foco}), mueve el
algoritmo de resolución de ecuaciones más rápido de uso corriente, dentro
de todas las calculadoras y de todas las redes neuronales, y encuentra la
viga rectangular más resistente que se puede sacar de un tronco redondo.

\medskip
\noindent\textbf{Parte I --- Soltura con la \omterm{def:g11:deriv:tangent}{tangente}.}
\begin{enumerate}
  \item Deriva $f(x) = 3x^2 - 5x + 1$ y $g(x) = x^3 - 12x$
        (\cref{thm:g11:deriv:rules}), y recupera $\left(x^2\right)'$ en
        $x = 1$ a partir de la definición (la tasa de variación y después
        $h \to 0$).
  \item Halla la \omterm{def:g11:deriv:tangent}{tangente} a $y = x^3 - 12x$ en $a = 2$. ¿Qué tiene de
        especial y qué señala eso (\cref{prop:g11:deriv:extrema})?
  \item Establece la \omterm{def:g11:deriv:tangent}{tangente} general a la \omterm{def:g11:quad:parabola}{parábola} $y = x^2$ en el punto
        $(a, a^2)$:
        \[
        y = 2ax - a^2 .
        \]
  \item ¿Dónde corta esa \omterm{def:g11:deriv:tangent}{tangente} al eje $y$? Deduce la receta del
        delineante: para trazar la \omterm{def:g11:deriv:tangent}{tangente} en un punto $P$ de altura $h$
        de la \omterm{def:g11:quad:parabola}{parábola}, une $P$ con el punto del eje que está a
        profundidad\,\dots\ por debajo del origen. Enúnciala.
  \item Construye la \omterm{def:g10:functions:table}{tabla de variación} completa de $g(x) = x^3 - 12x$
        (signo de $g'$, extremos relativos y sus valores,
        \cref{met:g11:deriv:study}).
\end{enumerate}

\medskip
\noindent\textbf{Parte II --- El teorema del espejo.}
Sea $P(a, a^2)$ un punto de la \omterm{def:g11:quad:parabola}{parábola} $y = x^2$ (con $a \neq 0$), sea
$F\left(0, \frac14\right)$ el \omterm{pb:g11:quad:1}{foco} y sea $T$ el punto donde la \omterm{def:g11:deriv:tangent}{tangente} en
$P$ corta al eje $y$ (pregunta 4). Recuerda del \cref{pb:g11:quad:1} que
$PF = a^2 + \frac14$.
\begin{enumerate}[resume]
  \item Calcula la distancia $FT$.
  \item Concluye que el triángulo $FPT$ es isósceles en $F$. ¿Qué dos
        ángulos son, por tanto, iguales?
  \item El rayo que llega a $P$ paralelo al eje es paralelo a la recta
        $(TF)$ (los dos son verticales). Usando los ángulos alternos,
        demuestra que la \omterm{def:g11:deriv:tangent}{tangente} en $P$ forma \emph{ángulos iguales} con
        el rayo vertical entrante y con el segmento $[PF]$.
  \item Física: la luz se refleja en una curva como lo haría en su
        \omterm{def:g11:deriv:tangent}{tangente}, y sale con el mismo ángulo con el que llegó (la ley de
        la reflexión). Concluye el \emph{teorema del espejo}: todo rayo
        paralelo al eje se refleja pasando por el \omterm{pb:g11:quad:1}{foco} y, invirtiendo el
        tiempo, una bombilla situada en el \omterm{pb:g11:quad:1}{foco} emite un haz paralelo. La
        antena y el faro del \cref{pb:g11:quad:1} son ya teoremas.
  \item Comprueba numéricamente el triángulo isósceles en $a = 1$:
        calcula $T$, $FP$ y $FT$.
\end{enumerate}

\medskip
\noindent\textbf{Parte III --- La máquina de Newton.}
Para resolver $f(x) = 0$, Newton propuso: partir de una estimación $x_0$,
seguir la \omterm{def:g11:deriv:tangent}{tangente} en $x_0$ hasta el eje y tomar su punto de corte como
estimación siguiente.
\begin{enumerate}[resume]
  \item Deduce la fórmula del método:
        \[
        x_{n + 1} = x_n - \frac{f(x_n)}{f'(x_n)} .
        \]
  \item Aplícalo a $f(x) = x^2 - 2$ desde $x_0 = 1$: calcula $x_1$, $x_2$
        y $x_3$ como fracciones exactas. ¿Qué sucesión de hace dos mil
        años, y qué máquina del \cref{pb:g10:functions:1}, acabas de
        reconstruir? Demuestra la identidad que hay detrás de esa
        coincidencia:
        \[
        x - \frac{x^2 - 2}{2x} = \frac12\left(x + \frac2x\right).
        \]
  \item Resuelve $x^3 = 100$ (el cruce de los planes de ahorro del
        \cref{pb:g10:reffunc:1}): aplica dos pasos de Newton a
        $f(x) = x^3 - 100$ desde $x_0 = 5$ (con cuatro decimales) y
        compara con el $\sqrt[3]{100}$ de la calculadora.
  \item Sabotaje: intenta arrancar el método con $f(x) = x^3 - 12x$ en
        $x_0 = 2$. ¿Qué falla y por qué (pregunta 2)? Enuncia la
        advertencia práctica.
  \item En la pregunta 12, los decimales correctos fueron más o menos
        $1 \to 3 \to 6$. Compara con la bisección (partir el \omterm{def:g10:numbers:interval}{intervalo} por
        la mitad en cada paso) y di en una frase por qué las calculadoras,
        los receptores de GPS y el entrenamiento de las redes neuronales
        funcionan todos con la idea de Newton.
\end{enumerate}

\medskip
\noindent\textbf{Parte IV --- La viga más resistente.}
De un tronco redondo de $30$ cm de diámetro se sierra una viga rectangular
de anchura $w$ y canto $d$, de modo que $w^2 + d^2 = 900$. La rigidez de
una viga es proporcional a $w d^2$ (el canto cuenta el doble: ¡las viguetas
se ponen de canto!).
\begin{enumerate}[resume]
  \item Expresa la rigidez en \omterm{def:g11:func:function}{función} de $w$ sola: $S(w) = 900w - w^3$.
        ¿En qué \omterm{def:g10:numbers:interval}{intervalo}?
  \item Deriva, construye la \omterm{def:g10:functions:table}{tabla de variación} y halla los valores
        óptimos de $w$ y $d$ (exactos y después al milímetro).
  \item Demuestra que las proporciones óptimas cumplen $d = w\sqrt2$: la
        vieja razón de los carpinteros. (Curiosidad adicional: dividir el
        diámetro del tronco en tres partes iguales y levantar
        perpendiculares hasta la circunferencia construye exactamente ese
        rectángulo.)
  \item La elección ingenua es la viga \emph{cuadrada} ($w = d$). Calcula
        su rigidez y la de la viga óptima, y da la ganancia del carpintero
        en porcentaje.
  \item Final: los tres oficios de la \omterm{def:g11:deriv:tangent}{tangente}, una frase para cada uno:
        la geometría (el teorema del espejo), el cálculo numérico (la
        máquina de Newton) y la optimización (la viga), todos movidos por
        el mismo objeto. Y el aviso hacia delante: el año que viene se
        añade la convexidad, o sea, de qué lado de la \omterm{def:g11:deriv:tangent}{tangente} está la
        curva, para afinar los tres.
\end{enumerate}
\end{problem}
